\clearpage
\fancyhead[L]{\small THEORY OF INTERPOLATION}
\pdfbookmark[0]{Theory of interpolation treated by a new method}{interpolation-nachlass}
\begingroup
\renewcommand{\arraystretch}{1.35}
\SourceLeaf{274}
\begin{center}\textbf{POSTHUMOUS PAPERS.}\par\bigskip
\textbf{THEORY OF INTERPOLATION}\par\medskip
\textbf{TREATED BY A NEW METHOD.}\par\medskip
\rule{.15\linewidth}{.4pt}\end{center}
\InterpolationArticle{1}
\textsc{Problem.}\quad\textit{Find the sum of the series}
\[
\begin{aligned}
&\frac{a^n}{(a-b)(a-c)(a-d)(a-e)\ldots}\\[5pt]
&+\frac{b^n}{(b-a)(b-c)(b-d)(b-e)\ldots}\\
&+\frac{c^n}{(c-a)(c-b)(c-d)(c-e)\ldots}\\[5pt]
&+\frac{d^n}{(d-a)(d-b)(d-c)(d-e)\ldots}\\
&+\frac{e^n}{(e-a)(e-b)(e-c)(e-d)\ldots}+\text{etc.}
\end{aligned}
\]
\textit{where $a,b,c,d,e$ are $m$ distinct quantities, and $n$ is any integer, positive, negative or even 0.}

\textit{Solution.}\quad For brevity, put
\[
\begin{aligned}
\frac1{(a-b)(a-c)(a-d)(a-e)\ldots}&=\alpha\\
\frac1{(b-a)(b-c)(b-d)(b-e)\ldots}&=\beta\\
\frac1{(c-a)(c-b)(c-d)(c-e)\ldots}&=\gamma\\
\frac1{(d-a)(d-b)(d-c)(d-e)\ldots}&=\delta,\ \text{etc.}
\end{aligned}
\]
so that the required sum, which we shall denote by $S^n$, becomes $=\alpha a^n+\beta b^n+\gamma c^n+\delta d^n+$ etc. It is evident that, if $x$ denotes an indeterminate quantity, expanding the sum
\SourceLeaf{275}
\[
P=\frac\alpha{1-ax}+\frac\beta{1-bx}+\frac\gamma{1-cx}+\frac\delta{1-dx}+\text{etc.}
\]
in a series of ascending powers of $x$ gives
\[
S^0+S^1x+S^2xx+S^3x^3+\text{etc.}\quad\text{to infinity}
\]
Put $(1-ax)(1-bx)(1-cx)(1-dx)\ldots=Q$; then $Q$ will be a polynomial in the indeterminate $x$, reaching the $m^{\mathrm{th}}$ degree; $PQ$, however, will be a polynomial of the $m-1^{\mathrm{th}}$ degree, namely $=$
\[
\begin{aligned}
&\alpha(1-bx)(1-cx)(1-dx)\ldots\\
&+\beta(1-ax)(1-cx)(1-dx)\ldots\\
&+\gamma(1-ax)(1-bx)(1-dx)\ldots\\
&+\delta(1-ax)(1-bx)(1-cx)\ldots\\
&+\text{etc.}
\end{aligned}
\]
A closer examination shows that the substitution $x=\frac1a$ makes all terms except the first vanish, while the first becomes
\[
\alpha(1-\tfrac ba)(1-\tfrac ca)(1-\tfrac da)\ldots=\frac1{a^{m-1}}
\]
In the same way, the substitution $x=\frac1b$ makes all terms except the second vanish, and this becomes $=\frac1{b^{m-1}}$. Likewise, the substitutions $x=\frac1c$, $x=\frac1d$ etc. turn $PQ$ into $\frac1{c^{m-1}}$, $\frac1{d^{m-1}}$ etc. It follows that $PQ-x^{m-1}$ takes the value 0 under all these substitutions, which is impossible unless it is identically $=0$, or $PQ=x^{m-1}$; otherwise the equation $PQ-x^{m-1}=0$, whose degree is no higher than the $m-1^{\mathrm{th}}$, would have $m$ distinct roots $\frac1a,\frac1b,\frac1c,\frac1d\ldots$.

Now let
\[
Q=1-Ax+Bxx-Cx^3+Dx^4-\text{etc.}
\]
where $A$ is the sum of the quantities $a,b,c,d\ldots$; $B$ the sum of their products taken two at a time; $C$ the sum of their products taken three at a time, etc. Since expanding the fraction $\frac{x^{m-1}}Q=P$ must give
\[
S^0+S^1x+S^2xx+S^3x^3+\text{etc.}
\]
it is evident, first, that the following must hold: $S^0=0$, $S^1=0$, $S^2=0$ etc., up to $S^{m-2}=0$\HistoricalNote{22}; then
\SourceLeaf{276}
the relations $S^{m-1}=1$, $S^m=A$ hold, and finally the later terms are determined, as terms of a recurrent series, by the following rule:
\[
\begin{aligned}
S^m&=A\\
S^{m+1}&=AS^m-BS^{m-1}=AA-B\\
S^{m+2}&=AS^{m+1}-BS^m+CS^{m-1}=A^3-2BA+C\\
&\text{etc.}
\end{aligned}
\]
\Needspace{14\baselineskip}
It is indeed easy to infer from this that $S^{m+1}$ $=aa+bb+cc+\ldots+B$, or the sum of the squares together with the sum of all products of two distinct quantities among $a,b,c,d\ldots$. But to see more clearly how the later terms are formed from the elements $a,b,c,d\ldots$, we observe that $\frac1Q$ is the product of the series
\[
\begin{aligned}
&1+ax+aaxx+a^3x^3+\text{etc.}\\
&1+bx+bbxx+b^3x^3+\text{etc.}\\
&1+cx+ccxx+c^3x^3+\text{etc.}\\
&1+dx+ddxx+d^3x^3+\text{etc.}\\
&\text{etc.}
\end{aligned}
\]
This product is $=\Sigma a^\lambda b^\mu c^\nu\ldots\times x^{\lambda+\mu+\nu.\,.}$, where the exponents $\lambda,\mu,\nu\ldots$ must be assigned all integer values from 0 to infinity and combined in every possible way. Thus, in the series into which $\frac1Q$ expands, to obtain the coefficient of the term containing $x^{n+1-m}$ we must split the number $n+1-m$ in every possible way into $n$ integer parts $\lambda+\mu+\nu+.\,.$ (among which a part 0 is also allowed), admitting every permutation of these parts as well; we then add the individual products $a^\lambda b^\mu c^\nu\ldots$. Their sum will be the required coefficient, and at the same time $=S^n$. A little attention shows that the following rule is entirely equivalent: from the $m$ quantities $a,b,c,d\ldots$ take all combinations of $n+1-m$ elements, allowing repetitions, and regard each combination as a product; the sum of these products will be $=S^n$. Accordingly, as above, $S^{m+1}$ will be the sum of all products of two quantities among $a,b,c,d\ldots$, whether distinct or identical; $S^{m+2}$ the sum of all products of three, distinct or identical, etc.

It only remains for us to show how to determine the sum of our progression for negative values of $n$. For this purpose, we shall put the first term of the sum $S^{-n}$, namely $\frac{a^{-n}}{(a-b)(a-c)(a-d)\ldots}$, in the following form
\SourceLeaf{277}
\[
\pm\frac1{abcd\ldots}\times\frac{(\frac1a)^{m+n-2}}{(\frac1a-\frac1b)(\frac1a-\frac1c)(\frac1a-\frac1d)\ldots}
\]
\Needspace{15\baselineskip}
where the upper or lower sign is to be taken according as $m$ is odd or even; the same transformation can also be applied to the remaining terms. Therefore, if the symbol $T$ denotes the expression formed from $\frac1a,\frac1b,\frac1c,\frac1d.\,.$ in the same way as $S$ is formed from $a,b,c,d\ldots$, evidently $S^{-n}=\pm\frac{T^{m+n-2}}{abcd\ldots}$. This case has thus been reduced to the preceding one, and the result is
\[
S^{-1}=\pm\frac1{abcd\ldots},\qquad S^{-2}=\pm\frac{\frac1a+\frac1b+\frac1c+\frac1d+.\,.}{abcd\ldots},
\]
$S^{-3}$ equals the product of $\pm\frac1{abcd\ldots}$ and the sum of all products of two quantities among $\frac1a,\frac1b,\frac1c,\frac1d\ldots$, distinct or identical, etc.
\InterpolationArticle{2}
We shall apply the preceding investigation to the case in which the quantities $a,b,c,d$ are given imaginary values. In this way we shall very easily obtain some remarkable relations that would present greater difficulties if treated by another method. Let $E$ be the base of natural logarithms, and $i$ the imaginary quantity $\BareRoot{-1}$; in place of the real quantities $a,b,c\ldots$, let us consider the imaginary quantities $E^{ia},E^{ib},E^{ic},E^{id}\ldots$ and $E^{-ia},E^{-ib},E^{-ic},E^{-id}\ldots$, and put
\[
\begin{aligned}
&\frac{E^{ina}}{(E^{ia}-E^{ib})(E^{ia}-E^{ic})(E^{ia}-E^{id})\ldots}\\
&+\frac{E^{inb}}{(E^{ib}-E^{ia})(E^{ib}-E^{ic})(E^{ib}-E^{id})\ldots}\\
&+\frac{E^{inc}}{(E^{ic}-E^{ia})(E^{ic}-E^{ib})(E^{ic}-E^{id})\ldots}\\
&+\frac{E^{ind}}{(E^{id}-E^{ia})(E^{id}-E^{ib})(E^{id}-E^{ic})\ldots}+\text{etc.}=S^n,\quad\text{and}
\end{aligned}
\]
\[
\begin{aligned}
&\frac{E^{-ina}}{(E^{-ia}-E^{-ib})(E^{-ia}-E^{-ic})(E^{-ia}-E^{-id})\ldots}\\
&+\frac{E^{-inb}}{(E^{-ib}-E^{-ia})(E^{-ib}-E^{-ic})(E^{-ib}-E^{-id})\ldots}\\
&+\frac{E^{-inc}}{(E^{-ic}-E^{-ia})(E^{-ic}-E^{-ib})(E^{-ic}-E^{-id})\ldots}\\
&+\frac{E^{-ind}}{(E^{-id}-E^{-ia})(E^{-id}-E^{-ib})(E^{-id}-E^{-ic})\ldots}+\text{etc.}=T^n
\end{aligned}
\]
\SourceLeaf{278}
\Needspace{14\baselineskip}
Thus, if the number of quantities $a,b,c,d\ldots$ is denoted by $m$, the sums $S^0,S^1,S^2\ldots S^{m-2}$, as well as $T^0,T^1,T^2\ldots T^{m-2}$, will all be $=0$; further, $S^{m-1}=T^{m-1}=1$; $S^m$ will be the sum of the quantities $E^{ia},E^{ib},E^{ic},E^{id}\ldots$; $S^{m+1}$ the sum of all their products taken two at a time, distinct or identical; $S^{m+2}$ the sum of their products taken three at a time, etc. In the same way, the values of the sums $T^m,T^{m+1},T^{m+2}$ etc. are to be formed from the quantities $E^{-ia},E^{-ib},E^{-ic},E^{-id}\ldots$.

\Needspace{21\baselineskip}
Now, since it is known that $E^{ix}+E^{-ix}=2\cos x$, $E^{ix}-E^{-ix}=2i\sin x$, it is easy to see that the value of the expression $\frac12(S^n+T^n)$, which for $n=0,1,2\ldots m-2$ is $=0$ and for $n=m-1$ is $=1$, becomes, for $n=m$,
\[
=\cos a+\cos b+\cos c+\cos d\,.\,.\,.\,.
\]
and likewise that $\frac12(S^{m+1}+T^{m+1})$ is the sum of the cosines of all angles obtained by adding \textit{two} of the angles $a,b,c,d.\,.$, distinct or identical; $\frac12(S^{m+2}+T^{m+2})$ is the sum of the cosines of all angles obtained by adding \textit{three} of the same angles, etc. Similarly, $\frac{S^n-T^n}{2i}=0$ for $n=0,1,2\ldots m-1$; further, it is $=\sin a+\sin b+\sin c+\sin d+.\,.$ for $n=m$; and likewise $\frac{S^{m+1}-T^{m+1}}{2i}$ will be the sum of the sines of all angles obtained by adding two of the angles $a,b,c,d\ldots$, distinct or identical; $\frac{S^{m+2}-T^{m+2}}{2i}$ the sum of the sines of all angles obtained by combining three of the same angles, etc.

We shall now examine the terms of the sums $S^n,T^n$ more closely. We have
\[
\begin{aligned}E^{ia}-E^{ib}&=E^{\frac12i(a+b)}(E^{\frac12i(a-b)}-E^{\frac12i(b-a)})\\
&=2iE^{\frac12i(a+b)}\sin\tfrac12(a-b)\end{aligned}
\]
and likewise
\[
E^{ia}-E^{ic}=2iE^{\frac12i(a+c)}\sin\tfrac12(a-c)\quad\text{etc.}
\]
Hence the denominator of the first term in $S^n$ becomes
\[
=(2i)^{m-1}E^{\frac12i(b+c+d\ldots+(m-1)a)}\sin\tfrac12(a-b)\sin\tfrac12(a-c)\sin\tfrac12(a-d)\ldots
\]
and, putting $a+b+c+d+\ldots=s$, this term itself becomes
\[
=\frac{E^{i((n+1-\frac12m)a-\frac12s)}}{(2i)^{m-1}\sin\tfrac12(a-b)\sin\tfrac12(a-c)\sin\tfrac12(a-d)\ldots}
\]
In the same way, the following holds:
\[
E^{-ia}-E^{-ib}=-2iE^{-\frac12i(a+b)}\sin\tfrac12(a-b)
\]
from which the first term of the sum $T^n$ finally becomes
\SourceLeaf{279}
\[
=\frac{\pm E^{-i((n+1-\frac12m)a-\frac12s)}}{(2i)^{m-1}\sin\tfrac12(a-b)\sin\tfrac12(a-c)\sin\tfrac12(a-d)\ldots}
\]
\Needspace{24\baselineskip}
where the upper or lower sign applies according as $m$ is odd or even. Adding this to the first term of the sum $S^n$, or subtracting it from that term, we conclude, \textit{first}, that for odd $m$ the first term of the sum $\frac12(S^n+T^n)$ becomes
\[
=\frac{\cos((n+1-\frac12m)a-\frac12s)}{(2i)^{m-1}\sin\tfrac12(a-b)\sin\tfrac12(a-c)\sin\tfrac12(a-d)\ldots}
\]
and the first term of the sum $\frac{S^n-T^n}{2i}$ becomes
\[
=\frac{\sin((n+1-\frac12m)a-\frac12s)}{(2i)^{m-1}\sin\tfrac12(a-b)\sin\tfrac12(a-c)\sin\tfrac12(a-d)\ldots}
\]
\textit{second}, that for even $m$ the first term of the sum $\frac12(S^n+T^n)$ becomes
\[
=\frac{\sin((n+1-\frac12m)a-\frac12s)}{2^{m-1}i^{m-2}\sin\tfrac12(a-b)\sin\tfrac12(a-c)\sin\tfrac12(a-d)\ldots}
\]
and the first term of the sum $\frac{S^n-T^n}{2i}$ becomes
\[
=\frac{-\cos((n+1-\frac12m)a-\frac12s)}{2^{m-1}i^{m-2}\sin\tfrac12(a-b)\sin\tfrac12(a-c)\sin\tfrac12(a-d)\ldots}
\]
The later terms of the expressions $\frac12(S^n+T^n)$, $\frac{S^n-T^n}{2i}$ will evidently be entirely analogous to the first, and will arise from it simply by interchanging the symbol $a$ with $b,c,d\ldots$. Now denote an arbitrary angle by $k$ and put $n+1-\frac12m=\lambda$. With the aid of the formulas
\[
\begin{aligned}
\cos(\lambda a+k)&=\cos(k+\tfrac12s)\cos(\lambda a-\tfrac12s)\\&\qquad-\sin(k+\tfrac12s)\sin(\lambda a-\tfrac12s)\\
\sin(\lambda a+k)&=\cos(k+\tfrac12s)\sin(\lambda a-\tfrac12s)\\&\qquad+\sin(k+\tfrac12s)\cos(\lambda a-\tfrac12s)
\end{aligned}
\]
\Needspace{24\baselineskip}
we shall without difficulty obtain the sums of the following series
\[
\begin{aligned}
&\frac{\cos(\lambda a+k)}{\sin\tfrac12(a-b)\sin\tfrac12(a-c)\sin\tfrac12(a-d)\ldots}\\
&+\frac{\cos(\lambda b+k)}{\sin\tfrac12(b-a)\sin\tfrac12(b-c)\sin\tfrac12(b-d)\ldots}\\
&+\frac{\cos(\lambda c+k)}{\sin\tfrac12(c-a)\sin\tfrac12(c-b)\sin\tfrac12(c-d)\ldots}\\
&+\frac{\cos(\lambda d+k)}{\sin\tfrac12(d-a)\sin\tfrac12(d-b)\sin\tfrac12(d-c)\ldots}+\text{etc.}=U^\lambda
\end{aligned}
\]
and
\SourceLeaf{280}
\[
\begin{aligned}
&\frac{\sin(\lambda a+k)}{\sin\tfrac12(a-b)\sin\tfrac12(a-c)\sin\tfrac12(a-d)\ldots}\\
&+\frac{\sin(\lambda b+k)}{\sin\tfrac12(b-a)\sin\tfrac12(b-c)\sin\tfrac12(b-d)\ldots}\\
&+\frac{\sin(\lambda c+k)}{\sin\tfrac12(c-a)\sin\tfrac12(c-b)\sin\tfrac12(c-d)\ldots}\\
&+\frac{\sin(\lambda d+k)}{\sin\tfrac12(d-a)\sin\tfrac12(d-b)\sin\tfrac12(d-c)\ldots}+\text{etc.}=V^\lambda.
\end{aligned}
\]
\Needspace{27\baselineskip}
\begin{center}\textit{First case: $m$ is odd.}\end{center}
In this case $\lambda$ must be, not an integer, but an integer with the fraction $\frac12$ added, or half an odd number. There is no need to consider \textit{negative} values of $\lambda$, since $\cos(-\lambda a+k)=\cos(\lambda a-k)$ and $\sin(-\lambda a+k)=-\sin(\lambda a-k)$; thus the sum for a negative value follows directly from that for the opposite positive value simply by changing $k$ into $-k$. This observation will evidently apply in the next case too. It is now easy to see that for $n=0,1,\ldots m-2$, or for $\lambda=-\frac12m+1,-\frac12m+2,\ldots\frac12m-1$, or, if we omit the negative values, for $\lambda=\frac12,\frac32,\frac52\ldots\frac12m-1$, both $U^\lambda=0$ and $V^\lambda=0$ hold; further,
\[
\begin{aligned}
U^{\frac12m}&=(2i)^{m-1}\cos(\tfrac12s+k),\\V^{\frac12m}&=(2i)^{m-1}\sin(\tfrac12s+k)\\
U^{\frac12m+1}&=(2i)^{m-1}\{\cos(\tfrac12s+k+a)+\cos(\tfrac12s+k+b)\\&\qquad+\cos(\tfrac12s+k+c)\\
&\qquad+\cos(\tfrac12s+k+d)+\text{etc.}\}\\
V^{\frac12m+1}&=(2i)^{m-1}\{\sin(\tfrac12s+k+a)+\sin(\tfrac12s+k+b)\\&\qquad+\sin(\tfrac12s+k+c)\\
&\qquad+\sin(\tfrac12s+k+d)+\text{etc.}\}
\end{aligned}
\]
$U^{\frac12m+2}$ or $V^{\frac12m+2}$ equals the product of $(2i)^{m-1}$ and the sum of the cosines or sines of all angles obtained by adding $\frac12s+k$ to two of the angles $a,b,c,d\ldots$, distinct or identical; $U^{\frac12m+3}$ or $V^{\frac12m+3}$ equals the product of $(2i)^{m-1}$ and the sum of the cosines or sines of all angles obtained by adding $\frac12s+k$ to three of the same angles, etc. Moreover, there is scarcely any need to point out that a power of the imaginary quantity $i$ is $=+1$ when its exponent is twice an even number, and $=-1$ when it is twice an odd number.
\Needspace{24\baselineskip}
\begin{center}\textit{Second case: $m$ is even.}\end{center}
In this case $\lambda$ must be an integer, and for $\lambda=0,1,2\ldots\frac12m-1$ both $U^\lambda=0$ and $V^\lambda=0$ hold. Further,
\SourceLeaf{281}
\[
\begin{aligned}
U^{\frac12m}&=2^{m-1}i^{m-2}\sin(\tfrac12s+k),\\
V^{\frac12m}&=-2^{m-1}i^{m-2}\cos(\tfrac12s+k);\\
U^{\frac12m+1}&=2^{m-1}i^{m-2}\{\sin(\tfrac12s+k+a)+\sin(\tfrac12s+k+b)\\&\qquad+\sin(\tfrac12s+k+c)\\
&\qquad+\sin(\tfrac12s+k+d)+\text{etc.}\}\\
V^{\frac12m+1}&=-2^{m-1}i^{m-2}\{\cos(\tfrac12s+k+a)+\cos(\tfrac12s+k+b)\\&\qquad+\cos(\tfrac12s+k+c)\\
&\qquad+\cos(\tfrac12s+k+d)+\text{etc.}\}
\end{aligned}
\]
$U^{\frac12m+2},U^{\frac12m+3}$ etc. equal the product of $2^{m-1}i^{m-2}$ and the sum of the sines of all angles obtained by adding $\frac12s+k$ to two, three, etc. of the angles $a,b,c\ldots$, distinct or identical; $V^{\frac12m+2},U^{\frac12m+3}$ etc. equal the product of $-2^{m-1}i^{m-2}$ and the sum of the cosines of all angles obtained by adding $\frac12s+k$ to two, three, etc. of the same angles.
\endgroup
